\documentclass[../main.tex]{subfiles}
\begin{document}

% capitulos/04_disenyo_experimental.tex

\chapter{Diseño experimental y metodología}
\label{C:disenyo}

\section{Introducción al diseño experimental}

El diseño adoptado es un \textbf{diseño factorial completo} $4 \times 3 \times 10$:
cuatro escenarios de ataque, tres estrategias de \emph{backup} y diez repeticiones
por celda, totalizando 120~ejecuciones completamente automatizadas. Este tipo de
diseño permite estudiar no solo el efecto principal de cada factor sino también
las posibles interacciones entre ellos: por ejemplo, si la ventaja de la estrategia~C
sobre la estrategia~B es diferente según el escenario de ataque.

La elección de diez repeticiones por celda responde a un equilibrio entre la
potencia estadística deseada y el tiempo de experimentación disponible. Con
$n = 10$ por celda es posible estimar la mediana con intervalos de confianza
razonables mediante \emph{bootstrap} no paramétrico, y ejecutar las pruebas de
Mann-Whitney~U con potencia suficiente para detectar efectos de tamaño mediano o
grande. Para detectar efectos pequeños (delta de Cliff $< 0.33$) con potencia del
80\,\% se requerirían al menos 30 observaciones por celda, lo que supondría
360~ejecuciones y un tiempo estimado de 15-20 horas. Esta limitación se documenta
en la sección de trabajo futuro.

\section{Variables del experimento}

El experimento sigue un diseño factorial $4 \times 3 \times 10$ con las siguientes
variables:

\subsection{Variable independiente principal: estrategia de backup}

\begin{description}
  \item[\stratA\ (control negativo)] Sin copia de seguridad. Representa el peor
  caso posible y permite cuantificar el impacto absoluto del ransomware sin
  mecanismo de recuperación. La única vía de recuperación en este caso sería el
  pago del rescate o la aceptación de la pérdida total de los datos.

  \item[\stratB\ (backup convencional)] Repositorio \restic sobre servidor
  \gls{sftp}. Los datos se almacenan en el contenedor \emph{restic-sftp},
  protegidos únicamente por autenticación \gls{ssh} con clave pública. Esta
  estrategia es representativa de una implementación técnicamente correcta de la
  regla 3-2-1, pero vulnerable a la destrucción activa del repositorio por un
  atacante con las credenciales comprometidas.

  \item[\stratC\ (backup inmutable)] Repositorio \restic sobre MinIO con Object
  Lock COMPLIANCE durante~30~días. Los objetos almacenados no pueden ser
  modificados ni borrados por ningún usuario durante el periodo de retención,
  incluyendo el usuario raíz de MinIO. Esta estrategia representa el estado del
  arte en protección de \emph{backups} frente al ransomware en entornos
  \emph{on-premise}.
\end{description}

\subsection{Variable independiente secundaria: escenario de ataque}
\label{sec:variables-escenario}

\begin{description}
  \item[E1 — Cifrado local] El atacante transfiere un \emph{script} cifrador en
  Python (AES-GCM, 256~bits) al \emph{fileserver} mediante SCP y lo ejecuta
  sobre \texttt{/opt/labprueba}. El cifrador renombra cada fichero añadiendo la
  extensión \texttt{.encrypted} y elimina el original. Corresponde a la técnica
  T1486 (\emph{Data Encrypted for Impact}).

  \item[E2 — Destrucción de copias] Ejecuta E1 más el intento de borrado del
  repositorio \gls{sftp} mediante \gls{ssh} y del \emph{bucket} de MinIO mediante
  AWS CLI con \texttt{aws s3 rm s3://restic-c/ --recursive}. Corresponde a
  T1490 (\emph{Inhibit System Recovery}).

  \item[E3 — Movimiento lateral] El atacante accede al \emph{fileserver} mediante
  \gls{ssh}, localiza la clave privada \texttt{id\_ed25519} en
  \texttt{/root/.ssh/} y la utiliza para establecer una nueva conexión SSH a la
  \emph{workstation}. Corresponde a T1021.004 (\emph{Remote Services: SSH}).

  \item[E4 — Exfiltración] El atacante empaqueta los ficheros del \emph{fileserver}
  en un archivo \texttt{.tar.gz} en \texttt{/tmp} y los transfiere al contenedor
  atacante mediante SCP. Corresponde a T1041 (\emph{Exfiltration Over C2
  Channel}).
\end{description}

\subsection{Variable de control: repetición}

Diez repeticiones por celda factorial ($n = 10$). Entre cada repetición, el
\emph{script} \texttt{reset\_lab.sh} restaura el laboratorio al estado inicial:
regenera los 500~ficheros sintéticos, recalcula los \emph{hashes} \emph{baseline},
limpia el repositorio \gls{sftp} y, según la estrategia, crea un nuevo repositorio
\restic con un \emph{snapshot} de los 500~ficheros. Para la estrategia~C se
esperan 120~segundos adicionales tras el \emph{backup} para que los eventos
\gls{fim} se procesen y no contaminen el \emph{baseline} de detección.

\section{Variables dependientes}

Se miden las siguientes variables dependientes en cada ejecución:

\begin{itemize}
  \item $\tdetect$ (segundos): tiempo entre el lanzamiento del escenario
  ($\tlaunch$) y la primera alerta en el \emph{indexer} para el grupo de reglas
  del escenario:
  \begin{equation}
    \Delta T_{\text{detect}} = T_{\text{detect}} - T_{\text{launch}}
  \end{equation}
  Se obtiene mediante \emph{polling} al \emph{indexer} cada 5~s. El intervalo de
  muestreo introduce una imprecisión máxima de $\pm 5$~s, aceptable dado que los
  tiempos observados oscilan entre 6 y 42~s.

  \item $\mttr$ (segundos): tiempo desde $\tlaunch$ hasta la finalización
  verificada de la restauración completa de los datos:
  \begin{equation}
    \text{MTTR} = T_{\text{restore}} - T_{\text{launch}}
  \end{equation}
  Esta definición incluye tanto el tiempo de detección como el tiempo de ejecución
  del \emph{playbook} de recuperación, reflejando el tiempo total de interrupción
  de la actividad.

  \item \texttt{recovery\_status} (categórico): \emph{success} cuando la
  restauración completa todos los ficheros con integridad verificada;
  \emph{impossible} cuando no existe copia disponible o fue destruida;
  \emph{partial} cuando la restauración es incompleta.

  \item \texttt{integrity\_match} (booleano): resultado de la comparación entre
  los hashes SHA-256 de los ficheros restaurados y el \emph{baseline} generado
  antes del ataque.

  \item \texttt{alert\_count} (entero): número de alertas nuevas generadas en el
  \emph{indexer} desde el inicio del ataque.
\end{itemize}

\section{Hipótesis estadísticas}

La hipótesis H1 se formula estadísticamente como:
\begin{align}
  H_{1,0} &: \text{la distribución del \gls{mttr} de B es idéntica a la de C} \\
  H_{1,a} &: \text{el \gls{mttr} de C es estocásticamente menor que el de B}
\end{align}

La hipótesis H2 se evalúa descriptivamente comparando las medianas de $\tdetect$
con el valor de referencia de 200~días reportado por IBM X-Force~2024 para la
detección manual.

La hipótesis H3 se formula como:
\begin{align}
  H_{3,0} &: p_B = p_C \quad \text{(tasas de éxito iguales ante E2)} \\
  H_{3,a} &: p_B < p_C \quad \text{(estrategia C tiene mayor tasa de éxito)}
\end{align}

\section{Pruebas estadísticas}

La distribución del \gls{mttr} no puede asumir normalidad a priori dado el tamaño
de muestra reducido ($n = 10$). Se emplean pruebas no paramétricas:

\begin{enumerate}
  \item \textbf{Mann-Whitney U} (bilateral, $\alpha = 0.05$): para comparar las
  distribuciones de \gls{mttr} entre estrategias B y C dentro de cada escenario.
  No asume normalidad ni igualdad de varianzas.

  \item \textbf{Delta de Cliff}~\cite{cliff1993} ($\delta$): para cuantificar el
  tamaño del efecto en escala $[-1, 1]$ sin asumir normalidad ni homogeneidad de
  varianzas, a diferencia de medidas paramétricas como la $d$ de Cohen:
  \begin{equation}
    \delta = \frac{\#(x_i > y_j) - \#(x_i < y_j)}{n_1 \cdot n_2}
  \end{equation}
  Interpretación según los umbrales de Romano et al.~\cite{romano2006}:
  negligible $|\delta| < 0.147$; pequeño $< 0.33$; mediano $< 0.474$; grande
  $\geq 0.474$.

  \item \textbf{Bootstrap} de intervalo de confianza al 95\,\% para la diferencia
  de medianas: 10.000~remuestreos con semilla fija (42) para reproducibilidad.
  Si el intervalo de confianza no contiene el cero, la diferencia es
  estadísticamente significativa al nivel $\alpha = 0.05$.

  \item \textbf{Test exacto de Fisher}: para comparar las tasas de éxito de
  recuperación entre estrategias sobre la tabla de contingencia $2\times 2$.
  Se prefiere al test $\chi^2$ porque las frecuencias esperadas son pequeñas
  ($n = 10$ por celda).
\end{enumerate}

\section{Reglas de detección Wazuh}
\label{sec:reglas}

Se desarrollaron 13~reglas \emph{custom} con identificadores en el rango
100200--100231. Dado que \texttt{auditd} no está disponible en el \emph{kernel}
CachyOS (véase sección~\ref{C:arquitectura}), todas las reglas se basan en dos
fuentes de datos:

\begin{enumerate}
  \item \gls{fim} (\texttt{syscheck}) con \texttt{realtime="yes"} en
  \texttt{/opt/labprueba} y \texttt{/tmp}.
  \item Eventos \gls{json} estructurados en \texttt{runner.jsonl} procesados como
  \texttt{localfile}.
\end{enumerate}

La tabla~\ref{tab:reglas} resume las reglas implementadas.

\begin{table}[ht]
  \centering
  \caption{Reglas Wazuh \emph{custom} del laboratorio PFG.}
  \label{tab:reglas}
  \rowcolors{2}{tabla-par}{white}
  \begin{tabularx}{\textwidth}{C{1.5cm} C{2.2cm} X C{1.2cm}}
    \toprule
    \rowcolor{tabla-cabecera}
    \textcolor{white}{\textbf{ID}} &
    \textcolor{white}{\textbf{Grupo}} &
    \textcolor{white}{\textbf{Trigger}} &
    \textcolor{white}{\textbf{Nivel}} \\
    \midrule
    100200 & e1\_local   & FIM masivo en /opt/labprueba (15~ev./120\,s) & 10 \\
    100201 & e1\_local   & Fichero .sh/.py creado en /tmp               & 12 \\
    100202 & e1\_local   & Ficheros .encrypted/.locked (10~ev./60\,s)   & 11 \\
    100205 & e1\_local   & runner: scenario=E1, event=execute            & 11 \\
    100206 & e1\_local   & runner: scenario=E1, event=end                & 10 \\
    100211 & e2\_destroy & FIM borrado masivo en /srv/backups/           & 13 \\
    100212 & e2\_destroy & runner: immutability\_verified                & 14 \\
    100213 & e2\_destroy & runner: immutability\_FAILED                  & 15 \\
    100214 & e2\_destroy & runner: phase2\_sftp\_delete\_attempt         & 12 \\
    100220 & e3\_lateral & runner: step3\_lateral\_jump                  & 10 \\
    100221 & e3\_lateral & runner: step1\_ssh\_fileserver                &  9 \\
    100230 & e4\_exfil   & runner: step2\_exfil\_command                 & 11 \\
    100231 & e4\_exfil   & runner: step3\_exfil\_complete                & 12 \\
    \bottomrule
  \end{tabularx}
\end{table}

\section{Procedimiento de cada ejecución}
\label{sec:procedimiento}

Cada una de las 120~ejecuciones sigue el siguiente procedimiento automatizado,
gestionado por \texttt{experiment\_runner.py}:

\begin{enumerate}
  \item \textbf{Reset}: el \emph{script} \texttt{reset\_lab.sh} regenera
  500~ficheros sintéticos en \texttt{/opt/labprueba}, calcula sus \emph{hashes}
  \emph{baseline}, limpia el repositorio \gls{sftp} y, según la estrategia, crea
  un nuevo repositorio \restic con un \emph{snapshot} de los 500~ficheros.

  \item \textbf{Baseline}: el \emph{runner} consulta el \emph{indexer} para
  obtener el número actual de alertas del grupo de reglas del escenario activo.

  \item \textbf{Ataque}: se lanza el escenario mediante \texttt{docker exec}
  sobre el contenedor atacante. Se registra $\tlaunch$ en nanosegundos mediante
  el reloj de alta resolución del sistema.

  \item \textbf{Detección}: el \emph{runner} realiza \emph{polling} al
  \emph{indexer} cada 5~s hasta detectar un delta positivo respecto al
  \emph{baseline}. Cuando $\Delta N > 0$, se registra $\tdetect$.

  \item \textbf{Recuperación}: se ejecuta el \emph{playbook} correspondiente a
  la estrategia. El \emph{playbook} verifica la disponibilidad del repositorio,
  limpia el directorio víctima, restaura los ficheros y comprueba la integridad
  por \emph{hash} SHA-256.

  \item \textbf{Registro}: se escribe una fila en el fichero \texttt{raw.csv}
  con todas las métricas. Se espera un \emph{cooldown} de 30~s antes de la
  siguiente ejecución.
\end{enumerate}

\section{Control de errores en la automatización}
\label{sec:control-errores}

La ejecución desatendida de 120~repeticiones a lo largo de aproximadamente
cuatro horas exige que el orquestador \texttt{experiment\_runner.py} sea
resiliente frente a fallos transitorios sin comprometer la validez de los datos
recogidos. El diseño adopta tres mecanismos de control de errores
complementarios.

En primer lugar, cada una de las seis fases del procedimiento descrito en la
sección~\ref{sec:procedimiento} --- \emph{reset}, \emph{baseline}, ataque,
detección, recuperación y registro --- se ejecuta dentro de un bloque con
tiempo máximo de espera (\emph{timeout}) explícito y diferenciado: 120~s para
el \emph{reset} completo del entorno, 90~s para el \emph{polling} de
detección (18~ciclos de 5~s, con margen sobre el máximo de 41.9~s observado en
la fase de calibración) y 180~s para la fase de recuperación, calibrado sobre el
peor caso observado en la estrategia~C del escenario~E4. Si cualquier fase
excede su \emph{timeout}, la ejecución se marca explícitamente con
\texttt{recovery\_status~=~timeout} en lugar de forzar un valor de MTTR
artificial, evitando así contaminar el análisis estadístico del
capítulo~\ref{C:resultados} con observaciones censuradas tratadas como
válidas. Como se indica en la sección~\ref{sec:experimento-ejecutado} del
capítulo~\ref{C:resultados}, ninguna de las 120~ejecuciones activó este
mecanismo durante el experimento final.

En segundo lugar, el orquestador verifica, antes de iniciar cada repetición, que
los nueve contenedores del laboratorio reportan estado \emph{healthy} o
\emph{Active} (según corresponda a Docker o al agente Wazuh), reutilizando la
misma lógica de comprobación que el \emph{script} \texttt{verify\_lab.sh}
descrito en el anexo~C. Si algún contenedor no está operativo, la ejecución se
pausa y reintenta la comprobación cada 10~s durante un máximo de 60~s antes de
abortar con un mensaje de diagnóstico, en lugar de proceder con un entorno
potencialmente inconsistente.

En tercer lugar, todos los errores no capturados por los dos mecanismos
anteriores --- por ejemplo, una excepción de red al consultar la \gls{api} del
\emph{indexer} --- se registran con su traza completa en un fichero de
\emph{log} separado del CSV de resultados (\texttt{runner\_errors.log}), y la
repetición afectada se repite automáticamente una única vez antes de marcarse
como fallida, siguiendo el principio de que un fallo de infraestructura
transitorio no debe descartarse como si fuera un resultado experimental
válido. Esta separación entre errores de infraestructura y resultados
experimentales es la razón por la que, como se señala en la
sección~\ref{sec:experimento-ejecutado} del capítulo~\ref{C:resultados}, el
CSV final no contiene ninguna fila con estados inconsistentes: cualquier
ejecución que no pudo completarse correctamente fue reintentada antes de
consolidar los resultados, en lugar de aparecer como una observación anómala en
el conjunto de datos analizado.

\section{Validez del diseño experimental}

\subsection{Validez interna}

La validez interna se garantiza mediante varios mecanismos. El \emph{reset}
completo entre ejecuciones asegura que cada repetición parte exactamente del
mismo estado inicial, eliminando la influencia de estados residuales. El registro
automatizado por el \emph{script} orquestador elimina errores de anotación manual.
La verificación de integridad por \emph{hash} SHA-256 garantiza que los ficheros
restaurados son byte a byte idénticos a los originales.

El orden de ejecución (A, B, C para cada escenario, en orden numérico de escenario)
es determinista, lo que facilita la reproducibilidad pero podría introducir sesgos
de orden si el estado del sistema cambia con el tiempo. Durante las cuatro horas
del experimento no se observó ninguna degradación del rendimiento que pudiera
atribuirse al orden de ejecución.

\subsection{Validez externa}

La validez externa está limitada por el uso de un entorno sintético. Los resultados
son directamente aplicables a organizaciones que usen arquitecturas de servidor de
ficheros Linux con \emph{backup} gestionado por \restic sobre \gls{sftp} o S3, y
donde el atacante tenga acceso \gls{ssh} con clave privada.

Los resultados son menos directamente aplicables a entornos con mayor
heterogeneidad tecnológica, sistemas Windows, o configuraciones de \emph{backup}
basadas en otras tecnologías. La extrapolación a cifradores reales requiere cautela
dado que el cifrador sintético utilizado no replica la velocidad de cifradores
como LockBit~3.0.

\subsection{Validez de constructo}
\label{sec:validez-constructo}

La validez de constructo evalúa si las métricas operacionalizadas
($\tdetect$, \gls{mttr},
\mbox{\texttt{recovery\_status}})
capturan efectivamente los conceptos teóricos de interés (capacidad de detección,
capacidad de recuperación).

El principal riesgo es que $\tdetect$, medido como tiempo hasta la primera alerta
del \emph{indexer}, no capture exactamente el instante en que el ataque comienza
a ser peligroso para la organización. En un entorno real, el atacante puede
permanecer en la red durante semanas antes de desplegar el cifrador, y la primera
alerta del \gls{siem} puede no coincidir con el inicio de la fase peligrosa del
ataque. En el laboratorio, $\tlaunch$ coincide exactamente con el inicio del
escenario, lo que elimina esta ambigüedad.

El \gls{mttr}, definido como tiempo desde el lanzamiento del ataque hasta la
verificación exitosa de la restauración, es comparable al \gls{rto} de los planes
de continuidad de negocio, aunque en un entorno real incluiría también el tiempo
de detección, notificación y toma de decisiones humanas que no se captura en el
laboratorio.

La variable \mbox{\texttt{recovery\_status}} captura correctamente la diferencia
cualitativa más importante del experimento: la distinción entre recuperación
exitosa, parcial e imposible.

\subsection{Validez de conclusión estadística}
\label{sec:validez-estadistica}

La validez de conclusión estadística evalúa si las relaciones observadas entre
las variables independientes y dependientes están correctamente respaldadas por
el análisis estadístico aplicado, con independencia de si dicha relación tiene
sentido teórico (validez de constructo) o es generalizable (validez externa).

Tres decisiones metodológicas de este trabajo están orientadas específicamente a
proteger esta dimensión de la validez. En primer lugar, el uso sistemático de
pruebas no paramétricas (Mann-Whitney~U, test exacto de Fisher) evita la
violación del supuesto de normalidad que se produciría al aplicar pruebas $t$ o
ANOVA sobre muestras de tamaño $n=10$, para las que el teorema central del
límite no ofrece garantías suficientes. En segundo lugar, la comunicación
conjunta del $p$-valor, el tamaño del efecto (delta de Cliff) y el intervalo de
confianza \emph{bootstrap} evita la interpretación errónea, frecuente en trabajos
con muestras pequeñas, de que una ausencia de significación estadística
($p > 0.05$) implica necesariamente ausencia de efecto real: el caso del
escenario~E4 (Cliff~$\delta = 0.360$, efecto mediano, $p = 0.186$) ejemplifica
precisamente esta distinción, y se discute en detalle en el
capítulo~\ref{C:discusion}. En tercer lugar, la semilla fija (42) empleada en el
remuestreo \emph{bootstrap} garantiza que los intervalos de confianza reportados
son exactamente reproducibles por terceros, evitando que la variabilidad del
propio procedimiento de remuestreo se confunda con variabilidad del fenómeno
estudiado.

El principal riesgo residual para la validez de conclusión estadística en este
trabajo es la reducida potencia estadística derivada de $n=10$ por celda, ya
señalada anteriormente en este capítulo: con este tamaño muestral, la
probabilidad de cometer un error de tipo~II --- no detectar un efecto que
realmente existe --- es considerable para efectos de tamaño pequeño o mediano,
como el observado en~E4. Este riesgo se documenta explícitamente como limitación
en el capítulo~\ref{C:discusion} y motiva la primera línea de trabajo futuro
propuesta en el capítulo~\ref{C:conclusiones}.

\end{document}
